\section{Analytic reconstruction of energy and refinement memory}
\label{sec:limite-energetico-k}

The positive separations constructed from the dodecaphasic register
determine an ordered partition and, by the preceding section, a
Dirichlet form on its nodal values. Its refinement has two complementary
operations: harmonic interpolation preserves the energy of inherited
data, whereas each detail records the variation added between those
data. Passage to unlimited depth consists in assembling this compatible
collection through a norm controlling the sum of its details. The result
concerns the one-dimensional analytic realization of energy; it retains
as its antecedent the HMT construction of the register and the joint
discrete structure of the continuum.

\subsection{Nested partitions and interpolation spaces}

Fix the initial partition
\(P_0=\{0=t_0<t_1<\cdots<t_{12}=1\}\), whose lengths are
\(d_j=t_j-t_{j-1}=K_j/\sum_iK_i>0\). Let
\((P_n)_{n\geq0}\) be a sequence of nested finite partitions of
\([0,1]\), obtained by subdividing the preceding intervals.
We write
\[
 |P_n|=\max\{b-a:[a,b]\text{ is an interval of }P_n\},
 \qquad \lim_{n\to\infty}|P_n|=0.
\]
Repeated subdivision of each interval into a fixed number of equal
successors, at least two, satisfies this condition. For other rules,
mesh contraction is an explicit hypothesis of this reader.
Successor labels and their carry data accompany the subdivision and
preserve the provenance of the refinement.

Let \(V_n\) be the complex space of continuous functions on \([0,1]\)
that are affine on each interval of \(P_n\). The energy and energy
norm are defined by
\begin{equation}
 \mathcal E_n(f)
 =\sum_{[a,b]\in P_n}\frac{|f(b)-f(a)|^2}{b-a}
 =\int_0^1|f'(x)|^2\,dx,
 \qquad
 \|f\|_{\mathcal E}^2=|f(0)|^2+\mathcal E_n(f).
 \label{eq:energia-afines-k}
\end{equation}
The integral equality follows because the derivative of an affine
function is constant on each interval. The energy thus agrees exactly
with the nodal form already constructed. The inclusion
\(V_n\subset V_{n+1}\) interprets the same affine function on a
finer partition; it preserves both \(f(0)\) and the energy and is
isometric for \(\|\cdot\|_{\mathcal E}\).

The initial value has a precise function: it fixes the constant
component identified by the Dirichlet form as its kernel. Energy
controls differences, and the pair consisting of the initial value
and those differences determines the observable. In the Sobolev space
\(H^1([0,1];\mathbb C)\), whose elements have an absolutely continuous
representative with weak derivative in \(L^2\), this norm is equivalent
to the usual norm. Indeed, \(f(x)=f(0)+\int_0^xf'(t)\,dt\)
controls the norm of \(f\) in terms of \(|f(0)|\) and
\(\|f'\|_{L^2}\); integration of the same identity in turn controls
\(|f(0)|\) through \(\|f\|_{L^2}+\|f'\|_{L^2}\).

\subsection{Reconstruction of a function from its compatible values}

\begin{theorem}[Energy reconstruction and completion of interpolations]
\label{thm:limite-energia-h1-k}
Let \(f_n\in V_n\) be a collection compatible under restriction to
inherited nodes:
\[
 f_{n+1}(a)=f_n(a)
 \quad\text{for every node }a\text{ of }P_n.
\]
If \(\sup_n\mathcal E_n(f_n)<\infty\), there is a unique function
\(f\in H^1([0,1];\mathbb C)\), considered in its continuous
representative, whose values at all these nodes are the prescribed
ones. Moreover, \(f_n\to f\) uniformly and in \(H^1\), and
\begin{equation}
 \int_0^1|f'(x)|^2\,dx
 =\lim_{n\to\infty}\mathcal E_n(f_n).
 \label{eq:limite-energia-h1-k}
\end{equation}
Conversely, nodal interpolation of any element of \(H^1\) produces
a collection of this kind. Thus the completion of \(\bigcup_n V_n\)
in the energy norm is \(H^1([0,1];\mathbb C)\).
\end{theorem}

\begin{proof}
Set \(g_n=f_n'\in L^2([0,1];\mathbb C)\). For \(m\geq n\),
compatibility of the values at the endpoints of each interval
\([a,b]\) of \(P_n\) gives
\[
 \frac1{b-a}\int_a^bg_m(x)\,dx
 =\frac{f_m(b)-f_m(a)}{b-a}
 =\frac{f_n(b)-f_n(a)}{b-a}
 =g_n\big|_{(a,b)}.
\]
Let \(W_n\) be the space of \(L^2\) functions constant on each
interval of \(P_n\), and let \(Q_n:L^2\to W_n\) be the orthogonal
projection given by averages over these intervals. The preceding
identity states \(Q_ng_m=g_n\). Consequently, \(g_m-g_n\) is
orthogonal to \(g_n\), and
\begin{equation}
 \|g_m-g_n\|_{L^2}^2
 =\mathcal E_m(f_m)-\mathcal E_n(f_n).
 \label{eq:pitagoras-gradientes-k}
\end{equation}
The energies are nondecreasing and bounded. The Pythagorean identity
proves that \((g_n)\) is a Cauchy sequence in \(L^2\), with
limit \(g\). Define
\[
 f(x)=f_0(0)+\int_0^xg(t)\,dt.
\]
This function belongs to \(H^1\), has weak derivative \(g\),
and, by the Cauchy--Schwarz inequality, satisfies
\begin{equation}
 \sup_{x\in[0,1]}|f_n(x)-f(x)|
 \leq\|g_n-g\|_{L^2}\longrightarrow0.
 \label{eq:convergencia-uniforme-energia-k}
\end{equation}
Uniform convergence preserves every prescribed nodal value; together
with convergence of the derivatives in \(L^2\), it gives convergence
in \(H^1\). Convergence of the derivative norms proves
\eqref{eq:limite-energia-h1-k}. The union of the nodes is dense
because \(|P_n|\to0\); two continuous representatives with the same
values on that union agree. This gives uniqueness.

For the converse, take \(f\in H^1\) in its continuous representative,
and let \(f_n\) be its interpolation at the nodes of \(P_n\).
Its derivative is \(g_n=Q_nf'\). The projections \(Q_n\) are
contractive on \(L^2\). For a continuous function \(v\), the
error \(|Q_nv-v|\) is bounded by its modulus of continuity evaluated
at \(|P_n|\), which tends to zero. Density of continuous functions
in \(L^2\) and contractivity imply \(Q_ng\to g\) for every
\(g\in L^2\): given a continuous approximant \(v\), use
\[
 \|Q_ng-g\|_{L^2}
 \leq 2\|g-v\|_{L^2}+\|Q_nv-v\|_{L^2}.
\]
Applying this convergence to \(g=f'\) gives convergence of the
derivatives and, through \eqref{eq:convergencia-uniforme-energia-k},
of the functions. Contractivity bounds their energies by
\(\|f'\|_{L^2}^2\), and convergence of the norms again proves
the limit identity. Every function in \(H^1\) is therefore an
energy limit of elements of \(\bigcup_nV_n\).
Equivalence of norms and completeness of \(H^1\) identify the
stated completion.
\end{proof}

The argument describes concretely what is preserved when passing from
the network to the function. Each finite derivative records the average
of all subsequent derivatives on its interval; refinement adds orthogonal
components that retain this average. Bounded energy controls the total
sum of these additions. The continuous value is then recovered by
integrating the limit derivative and preserving the initial value.
Through this explicit route, the analytic realization assembles the
compatible information of the finite levels.

\subsection{Orthogonal decomposition of memory by scales}

Let \(H_n\) be harmonic interpolation of the values on \(P_n\)
to the nodes of \(P_{n+1}\), interpreted as a function in
\(V_{n+1}\), and define the detail
\[
 z_{n+1}=f_{n+1}-H_nf_n.
\]
By \eqref{eq:detalle-energia-k}, the detail vanishes at inherited
nodes. Its derivative has zero average on each old interval and is
orthogonal to \(W_n\). For different levels, the derivative of an
earlier detail belongs to that space of constant functions; the
derivatives of the details are therefore mutually orthogonal.
The same orthogonality holds against \(f_0'\).

\begin{proposition}[Energy identity and recovery of the details]
\label{prop:memoria-energetica-escalas-k}
For every finite depth \(N\),
\begin{equation}
 \mathcal E_N(f_N)
 =\mathcal E_0(f_0)
 +\sum_{n=0}^{N-1}\mathcal E_{n+1}(z_{n+1}).
 \label{eq:energia-telescopica-detalles-k}
\end{equation}
If the collection satisfies the hypotheses of
Theorem~\ref{thm:limite-energia-h1-k}, the series on the right
converges and its sum is \(\int_0^1|f'|^2\). The data \(f_0\)
and all the details determine the complete compatible collection and
its analytic realization.
\end{proposition}
\begin{proof}
The finite identity follows by summing
\(\mathcal E_{n+1}(f_{n+1})
=\mathcal E_n(f_n)+\mathcal E_{n+1}(z_{n+1})\), the orthogonal
decomposition of the preceding section. The summands are nonnegative,
and the partial sums converge by \eqref{eq:limite-energia-h1-k}.
Finite reconstruction is performed recursively through
\(f_{n+1}=H_nf_n+z_{n+1}\); the preceding theorem then determines
the limit function. Conversely, a function in \(H^1\) determines
its interpolations and, by taking differences, all these details.
Each arrow thus preserves its own data.
\end{proof}

Variational elimination of interior nodes selects the zero detail
at the eliminated level. Preserving the details allows recovery of
any compatible finite-energy configuration. The effective response
and the complete configuration thus have a precise relationship:
the former records the minimum over the interior variables; the
latter adds the orthogonal components and their energy.

\subsection{Transport of the analytic realization under the positive flow}

The direction \(u=P_3P_{11}K\), recovered before energy evaluation
\cite{hmtX}, determines the separation collection in
\eqref{eq:schur-flujo-k}. Its parameter \(s\in\mathbb R\)
describes this analytic deformation. Write
\[
 d_j(s)=\frac{K_je^{su_j}}{\sum_\ell K_\ell e^{su_\ell}},
 \qquad
 t_i(s)=\sum_{j=1}^id_j(s),
 \qquad
 a_j(s)=\frac{d_j(s)}{d_j(0)}>0.
\]
Let \(F_s:[0,1]\to[0,1]\) be the map sending each \(t_i(0)\)
to \(t_i(s)\) and affine on the initial intervals. It is an
increasing homeomorphism, with slope \(a_j(s)\) on the interval
\([t_{j-1}(0),t_j(0)]\). Its inverse is also piecewise affine,
and both maps are Lipschitz.

Suppose each subdivision inherits the length fractions and direction
of its predecessor interval, as in the preceding section. All refined nodes
are then transported through the same \(F_s\). Denote their partitions
by \(P_n(s)=F_s(P_n(0))\) and their interpolation spaces by
\(V_n(s)\). The map
\[
 U_s f=f\circ F_s^{-1}
\]
transports \(V_n(0)\) to \(V_n(s)\), preserves corresponding
nodal values, and extends this action to \(H^1\).

\begin{theorem}[Energy equivalence and commutation with refinement]
\label{thm:transporte-h1-flujo-k}
For every \(f\in H^1([0,1];\mathbb C)\),
\begin{equation}
 \int_0^1|(U_sf)'(y)|^2\,dy
 =\sum_{j=1}^{12}\frac1{a_j(s)}
 \int_{t_{j-1}(0)}^{t_j(0)}|f'(x)|^2\,dx.
 \label{eq:transporte-energia-h1-k}
\end{equation}
For \(s\) in any compact interval, the energy norms before and
after transport are uniformly equivalent, and \(|P_n(s)|\to0\)
uniformly on that compact interval. Harmonic interpolation and
restriction to inherited nodes commute with \(U_s\); reconstruction
and minimization of the details are transported both at finite depth
and in the \(H^1\) realization.
\end{theorem}

\begin{proof}
On initial interval \(j\), the chain rule and change of variable
\(y=F_s(x)\) give \((U_sf)'(F_s(x))=f'(x)/a_j(s)\) and
\(dy=a_j(s)\,dx\). Integration and summation over the intervals
prove \eqref{eq:transporte-energia-h1-k}. The formula holds for
absolutely continuous representatives, and its right-hand side is finite.

Fix a compact parameter interval. Continuity and positivity of the
twelve factors give constants \(0<a_-\leq a_+<\infty\) such that
\(a_-\leq a_j(s)\leq a_+\) for every \(j\) and every \(s\)
in that compact interval. Therefore,
\[
 \frac1{a_+}\mathcal E(f)
 \leq\mathcal E(U_sf)
 \leq\frac1{a_-}\mathcal E(f),
 \qquad
 |P_n(s)|\leq a_+|P_n(0)|.
\]
Since \((U_sf)(0)=f(0)\), we also obtain
\[
 \min\{1,a_+^{-1}\}\|f\|_{\mathcal E}^2
 \leq\|U_sf\|_{\mathcal E}^2
 \leq\max\{1,a_-^{-1}\}\|f\|_{\mathcal E}^2.
\]
These bounds prove uniform equivalence and preservation of the
decreasing-mesh condition.

On each predecessor edge, the interpolation coefficients are the partial
sums of the inherited fractions, independent of \(s\).
If \(H_n(s)\) denotes its interpolation, then
\(U_sH_n(0)=H_n(s)U_s\) on the corresponding spaces.
Restriction commutes because the nodes are transported with their values.
The decomposition into harmonic function and detail is preserved,
and the Schur identity in \eqref{eq:schur-flujo-k} preserves its
variational characterization. Finally, bounded continuity of \(U_s\)
allows passage to the \(H^1\) limit. In particular, a collection
reconstructed before transport subsequently converges to \(U_sf\),
with the same transported nodal values.
\end{proof}

The law \eqref{eq:transporte-energia-h1-k} identifies the exact
metric factors of the deformation: the length change of each interval
inversely modifies its energy contribution. Equivalence of norms
preserves finite-energy configurations, and the detail archive preserves
their reconstruction. This gives a complete analytic realization of
this reader of the positive register. Its compatibility with canonical
forms is articulated on subdivisions of a fixed domain: summing forms
preserves the meromorphic object, whereas variational reduction and
the details preserve the energy response and refined observable.
Identification with higher-rank amplituhedral forms retains the
specific conditions of covering, degree, and residue correspondence
established for that realization.
